% !PRINTED 750
\weberpage{750}

Now it follows very simply that no number $g$ that coincides with one of the numbers $a, b$, or that lies between them when $a$ and $b$ are different, can belong to the sequence $\Omega$.

% [sic] -- the source prints singular ``kommt'' with the pair $\alpha_\nu,\beta_\nu$; English renders the collective meaning grammatically.
For the number $g$ lies in the interior of every one of the intervals $\delta_\nu$. Suppose that $g$ occurs in $\Omega$; if, following 3., we go so far with $\nu$ that the pair $\alpha_\nu, \beta_\nu$ comes later than $g$ in $\Omega$, then by 2. $g$ cannot lie in the interval $\delta_\nu$, and with that the impossibility of our assumption is proved.

\gesperrt{From this it follows, then, that the totality of the numbers of an interval $\alpha, \beta$ does not form an enumerable set.}

\begin{center}
\rule{4cm}{0.4pt}
\end{center}

This fact can be proved in yet another way, which is in a certain respect still simpler, and which we will set out in a few words. We do not restrict the generality if we take the interval from 0 to 1 as our basis. We think of all the numbers of this interval as represented by infinite decimal fractions. The finite decimal fractions are contained among them as well, if we set all figures, from a certain one on, equal to zero. In order to make the representation by decimal fractions a unique one, let it further be stipulated that for a finite decimal fraction this representation shall \gesperrt{always} be chosen, so that, e.\,g., 0,4999\dots\ shall not be put for 0,5000\dots.

We will now assume that these decimal fractions form an enumerable set; they can therefore be arranged in a countable sequence, which we represent thus:
\weq{$\Omega$}{
\begin{aligned}
\omega_1 &= 0,\ \alpha_1^{(1)}\, \alpha_2^{(1)}\, \alpha_3^{(1)}\, \dots\\
\omega_2 &= 0,\ \alpha_1^{(2)}\, \alpha_2^{(2)}\, \alpha_3^{(2)}\, \dots\\
\omega_3 &= 0,\ \alpha_1^{(3)}\, \alpha_2^{(3)}\, \alpha_3^{(3)}\, \dots\\
&\quad\cdots\cdots\cdots\cdots\cdots
\end{aligned}
}
where the $\alpha_\mu^{(\nu)}$ denote figures of the decadic system.

But it is now very easy to exhibit one decimal fraction (or indeed arbitrarily many) not contained in the sequence $\Omega$. We need only form
\[
\eta = 0,\ \beta_1\ \beta_2\ \beta_3\ \dots
\]
% German splits 'ver-|schieden' across the leaf but keeps $\alpha_\nu^{(\nu)}$ on this
% page. English cannot split 'dif-|ferent' and still carry 'from alpha' here, so the whole
% clause is absorbed onto this page; i0772 opens at the next sentence.
% [sic] -- the printed condition permits a recurring-9 expansion equivalent to an earlier terminating decimal; one repair is to choose $\beta_\nu\in\{1,2\}$ unequal to $\alpha_\nu^{(\nu)}$.
where the $\beta_\nu$ are figures of the decadic system that satisfy the single condition that $\beta_\nu$ shall for every $\nu$ differ from $\alpha_\nu^{(\nu)}$.
%CONT>>
